\section{Productos directos, extensiones y productos semidirectos}\label{sec:grupos-productos}

\subsection{Normales con intersección trivial}\label{tema:10-1-1}
Para subconjuntos $A,B\subseteq G$, escribimos
\[
[A,B]=\ang{[a,b]\given a\in A,\ b\in B}.
\]
\begin{lema}\label{lema:C10-conmutan}
Si $N,H\trianglelefteq G$, entonces $[N,H]\subseteq N\cap H$. En particular, si $N\cap H=\{e\}$, cada elemento de $N$ conmuta con cada elemento de $H$.
\end{lema}
\begin{proof}
Para $n\in N$ y $h\in H$, el elemento $[n,h]=n(hn^{-1}h^{-1})$ pertenece a $N$ por su normalidad. También $[n,h]=(nhn^{-1})h^{-1}\in H$. La intersección es subgrupo y contiene los generadores de $[N,H]$, luego contiene todo ese subgrupo. Si es trivial, cada conmutador es neutro y los elementos conmutan.
\end{proof}

\begin{prop}\label{prop:C10-directo}
Si $N,H\trianglelefteq G$ y $N\cap H=\{e\}$, la aplicación
\[
\mu:N\times H\longrightarrow NH,\qquad (n,h)\longmapsto nh
\]
es un isomorfismo. Si además $NH=G$, realiza $G$ como producto directo de esos subgrupos.
\end{prop}
\begin{proof}
El conjunto $NH$ es subgrupo por el \cref{lema:C8-cocientes}. El lema anterior da
\[
\mu(n_1,h_1)\mu(n_2,h_2)=n_1h_1n_2h_2
=n_1n_2h_1h_2=\mu(n_1n_2,h_1h_2).
\]
La aplicación es sobreyectiva por definición de $NH$. Si $nh=e$, entonces $n=h^{-1}\in N\cap H$, luego $n=h=e$. Su núcleo es trivial y es un isomorfismo.
\end{proof}

\begin{ejem}[Otro argumento para el orden $pq$]
Sean $q<p$ primos y $|G|=pq$. Si \mbox{$p\not\equiv1\pmod q$}, ambos subgrupos de Sylow son únicos: $n_p\mid q$ y $n_p\equiv1\pmod p$ fuerzan $n_p=1$, porque $q<p$; por otra parte, $n_q\mid p$ deja los candidatos $1,p$, y la hipótesis excluye $p$.

Sean $P,Q$ esos subgrupos normales. Por Lagrange, $P\cap Q=\{e\}$, pues su orden divide a $p$ y a $q$. La proposición da $PQ\cong P\times Q$, de orden $pq$, así que $PQ=G$. Como $P\cong C_p$ y $Q\cong C_q$, obtenemos $G\cong C_p\times C_q\cong C_{pq}$. La última identificación se justifica porque un par de generadores tiene orden $\mcm(p,q)=pq$.
\end{ejem}
\begin{obs}
Para $q<p$, la condición que excluye $n_q=p$ es $p\not\equiv1\pmod q$. La condición $q\not\equiv1\pmod p$ no basta: $S_3$, de orden $2\cdot3$, la satisface y no es cíclico.
\end{obs}

\begin{ejem}
En $S_3$, tomemos $N=\ang{(1\,2\,3)}$ y $H=\ang{(1\,2)}$. El primero es normal, por ser $A_3$; el segundo no lo es, ya que
\[
(1\,2\,3)(1\,2)(1\,2\,3)^{-1}=(2\,3)\notin H.
\]
Tenemos $N\cap H=\{e\}$ y $NH=S_3$: los seis productos $nh$ son distintos, pues $nh=n'h'$ implica $n'^{-1}n=h'h^{-1}\in N\cap H$. Sin embargo, $S_3\not\cong N\times H$. El producto directo tiene el elemento $((1\,2\,3),(1\,2))$ de orden seis, mientras que en $S_3$ los tipos posibles sólo dan órdenes uno, dos o tres. La normalidad de ambos factores es esencial para la proposición.
\end{ejem}

\subsection{Sucesiones exactas y escisiones}\label{tema:10-1-2}
\begin{defi}
Una sucesión de homomorfismos es \emph{exacta} en un término si la imagen de la flecha que llega es el núcleo de la que sale. Una \emph{extensión de $H$ por $N$} es una sucesión exacta
\[
\begin{tikzcd}[column sep=large]
1\arrow[r]&N\arrow[r,"i"]&G\arrow[r,"p"]&H\arrow[r]&1.
\end{tikzcd}
\]
Esto equivale a que $i$ sea inyectiva, $p$ sobreyectiva e $i(N)=\ker p$. Identificamos entonces $N$ con su imagen normal en $G$. La extensión \emph{se escinde} si existe un homomorfismo $s:H\to G$ tal que $p\circ s=\id_H$, llamado \emph{sección}.
\end{defi}
Toda sección es inyectiva: $s(h)=e$ implica $h=p(s(h))=e$. Su imagen es un subgrupo $H_0\le G$ y satisface
\[
N\cap H_0=\{e\},\qquad NH_0=G.
\]
En efecto, si $s(h)\in N$, entonces $h=p(s(h))=e$. Para cada $g\in G$, tomando $h=p(g)$, el elemento $gs(h)^{-1}$ está en $N$, así que $g\in NH_0$.

El producto directo siempre produce la sucesión exacta
\[
\begin{tikzcd}[column sep=large]
1\arrow[r]&N\arrow[r,"{n\mapsto(n,e)}"]&N\times H\arrow[r,"{(n,h)\mapsto h}"]&H\arrow[r]&1.
\end{tikzcd}
\]
La inclusión es inyectiva, la proyección es sobreyectiva y su núcleo es $N\times\{e\}$. La sección es $h\mapsto(e,h)$. Una sucesión exacta cualquiera no implica un producto directo: por ejemplo,
\[
1\longrightarrow A_3\longrightarrow S_3
\overset{\varepsilon}{\longrightarrow}\{1,-1\}\longrightarrow1
\]
es exacta y $S_3\not\cong C_3\times C_2$. Incluso se escinde, enviando $-1$ a $(1\,2)$; escindirse tampoco basta para ser producto directo.

\begin{lema}\label{lema:C10-escision}
Si $N\trianglelefteq G$, $H\le G$, $G=NH$ y $N\cap H=\{e\}$, entonces $G$ es una extensión escindida de $H$ por $N$.
\end{lema}
\begin{proof}
La restricción $\alpha:H\to G/N$, $h\mapsto hN$, es sobreyectiva: si $g=nh$, entonces $gN=h(h^{-1}nh)N=hN$, por normalidad de $N$. Su núcleo es $H\cap N=\{e\}$, por lo que es un isomorfismo. La aplicación $p=\alpha^{-1}\circ\pi:G\to H$ es sobreyectiva y tiene núcleo $N$. La inclusión $H\hookrightarrow G$ es una sección de $p$, de modo que la sucesión resultante es exacta y se escinde.
\end{proof}
\begin{obs}
Un complemento $H_0$ de $N$ debe satisfacer tanto $N\cap H_0=\{e\}$ como $NH_0=G$. La intersección trivial por sí sola no exige que ambos subgrupos generen $G$.
\end{obs}

\subsection{La operación determinada por una acción}\label{tema:10-2-1}
En la situación del lema anterior, la conjugación da un homomorfismo
\[
\gamma:H\longrightarrow\Aut(N),\qquad \gamma_h(n)=hnh^{-1}.
\]
La normalidad asegura que cada conjugación preserva $N$, su inversa es conjugar por $h^{-1}$, y $\gamma_{h_1h_2}=\gamma_{h_1}\circ\gamma_{h_2}$. Además,
\[
(n_1h_1)(n_2h_2)=\bigl(n_1\gamma_{h_1}(n_2)\bigr)(h_1h_2).
\]
Esto motiva la construcción abstracta siguiente.

\begin{prop}\label{prop:C10-construccion}
Sean $N,H$ grupos y $\theta:H\to\Aut(N)$ un homomorfismo. La operación en $N\times H$ dada por
\[
(n_1,h_1)\cdot_\theta(n_2,h_2)
=\bigl(n_1\theta_{h_1}(n_2),h_1h_2\bigr)
\]
define un grupo con neutro $(e_N,e_H)$ e inverso
\[
(n,h)^{-1}=\bigl(\theta_{h^{-1}}(n^{-1}),h^{-1}\bigr).
\]
\end{prop}
\begin{proof}
La operación está bien definida porque $\theta_h$ es un automorfismo de $N$. Para la asociatividad, el producto asociado a la izquierda es
\[
\bigl(n_1\theta_{h_1}(n_2)\theta_{h_1h_2}(n_3),h_1h_2h_3\bigr).
\]
El asociado a la derecha es
\[
\bigl(n_1\theta_{h_1}(n_2\theta_{h_2}(n_3)),h_1h_2h_3\bigr).
\]
Coinciden porque $\theta_{h_1}$ preserva productos y $\theta_{h_1h_2}=\theta_{h_1}\theta_{h_2}$. Como $\theta_{e_H}=\id_N$ y $\theta_h(e_N)=e_N$, multiplicar por $(e_N,e_H)$ a cualquiera de los dos lados deja el elemento intacto.

Para el inverso indicado, las primeras coordenadas de los productos en los dos órdenes son, respectivamente,
\[
n\theta_h(\theta_{h^{-1}}(n^{-1}))=nn^{-1}=e_N,
\qquad
\theta_{h^{-1}}(n^{-1})\theta_{h^{-1}}(n)=e_N.
\]
En ambos casos la segunda coordenada es $e_H$.
\end{proof}
\begin{defi}
El grupo de la proposición se llama \emph{producto semidirecto} y se denota $N\rtimes_\theta H$. Si $\theta$ es trivial, la operación es la del producto directo $N\times H$.
\end{defi}

\begin{prop}\label{prop:C10-propiedades}
En $G=N\rtimes_\theta H$, las aplicaciones $i(n)=(n,e_H)$ y $s(h)=(e_N,h)$ identifican copias de $N$ y $H$. La proyección $p(n,h)=h$ produce una extensión escindida, y esas copias cumplen $N\cap H=\{e_G\}$ y $NH=G$. La acción $\theta$ se realiza por conjugación en $G$.
\end{prop}
\begin{proof}
La fórmula de multiplicación muestra que $i$ y $s$ son homomorfismos inyectivos. También muestra que $p$ es un homomorfismo sobreyectivo con núcleo $i(N)$ y que $p\circ s=\id_H$. Así,
\[
\begin{tikzcd}[column sep=large]
1\arrow[r]&N\arrow[r,"i"]&N\rtimes_\theta H\arrow[r,"p"]&H\arrow[l,bend left=35,destacado,"s"]\arrow[r]&1
\end{tikzcd}
\]
expresa la sucesión exacta y su sección. Las dos copias se intersectan sólo en $(e_N,e_H)$ y todo $(n,h)$ es $(n,e_H)(e_N,h)$. Finalmente,
\begin{align*}
(e_N,h)(n,e_H)(e_N,h)^{-1}
&=(\theta_h(n),h)(e_N,h^{-1})\\
&=(\theta_h(n),e_H),
\end{align*}
que es la copia de $\theta_h(n)$ en $N$.
\end{proof}

\subsection{Producto interno y ejemplo de orden seis}\label{tema:10-2-2}
\begin{prop}\label{prop:C10-interno}
Si $N\trianglelefteq G$, $H\le G$, $G=NH$ y $N\cap H=\{e\}$, entonces $G\cong N\rtimes_\gamma H$, donde $\gamma$ es la acción por conjugación.
\end{prop}
\begin{proof}
Definimos $\varphi:N\rtimes_\gamma H\to G$ por $\varphi(n,h)=nh$. Es sobreyectiva por $G=NH$. Si $nh=n'h'$, entonces $n'^{-1}n=h'h^{-1}\in N\cap H$, así que $n=n'$ y $h=h'$; es biyectiva. Además,
\begin{align*}
\varphi\bigl((n_1,h_1)\cdot_\gamma(n_2,h_2)\bigr)
&=n_1(h_1n_2h_1^{-1})h_1h_2\\
&=(n_1h_1)(n_2h_2).
\end{align*}
Por tanto es un isomorfismo.
\end{proof}
Esta realización dentro de $G$ se llama \emph{producto semidirecto interno}. Si $N$ y $H$ conmutan, $\gamma$ es trivial y recuperamos el producto directo. En particular ocurre cuando ambos son normales y su intersección es trivial.

\begin{ejem}
Si $C_3=\ang{y}$, un automorfismo está determinado por la imagen de $y$, que debe ser $y$ o $y^2$. Así $\Aut(C_3)=\{\id,\iota\}\cong C_2$, donde $\iota(y)=y^2$ y $\iota^2=\id$. Hay exactamente dos homomorfismos $C_2\to\Aut(C_3)$: el trivial y el que envía su generador a $\iota$. El primero da $C_3\times C_2\cong C_6$; el segundo da un grupo isomorfo a $S_3$.
\end{ejem}
\begin{ejer}[Cálculo del producto semidirecto no trivial]
Comprobar que el producto semidirecto no trivial es $S_3$.
\end{ejer}
\begin{sol}
Escribamos $C_2=\ang{t}$ y elijamos $\theta_t(y)=y^{-1}$. Los seis elementos son $(y^a,t^b)$ con $a\in\{0,1,2\}$ y $b\in\{0,1\}$; su producto es
\[
(y^a,t^b)(y^c,t^d)=(y^{a+(-1)^bc},t^{b+d}),
\]
con exponentes reducidos módulo $3$ y $2$. En $S_3$ pongamos $r=(1\,2\,3)$ y $s=(1\,2)$. Por conjugación de ciclos, $srs^{-1}=(2\,1\,3)=r^{-1}$. En consecuencia,
\[
(r^as^b)(r^cs^d)=r^{a+(-1)^bc}s^{b+d}.
\]
La aplicación $(y^a,t^b)\mapsto r^as^b$ preserva los productos. Sus imágenes son $e,r,r^2,s,rs,r^2s$, todas distintas: los tres primeros elementos son pares y los otros impares, y dentro de cada bloque se cancela el mismo factor derecho. Por tanto es un isomorfismo sobre los seis elementos de $S_3$. También se obtiene aplicando la \cref{prop:C10-interno} a $N=A_3$ y $H=\ang{s}$.
\end{sol}
